% Generated draft: requires mathematical and source review.
\documentclass[11pt]{article}
\usepackage[margin=25mm]{geometry}
\usepackage{fontspec}
\setmainfont{Latin Modern Roman}
\usepackage{amsmath,amssymb,mathtools,unicode-math}
\setmathfont{Latin Modern Math}
\usepackage{enumitem,needspace,xurl,hyperref}
\setlength{\parindent}{0pt}
\setlength{\parskip}{5pt}
\setlength{\emergencystretch}{3em}
\newcommand{\N}{\mathbb{N}}
\newcommand{\Z}{\mathbb{Z}}
\newcommand{\Q}{\mathbb{Q}}
\newcommand{\R}{\mathbb{R}}
\newcommand{\C}{\mathbb{C}}
\newcommand{\F}{\mathbb{F}}
\newcommand{\A}{\mathbb{A}}
\newcommand{\PP}{\mathbb{P}}
\newcommand{\E}{\mathbb{E}}
\newcommand{\eps}{\varepsilon}
\newcommand{\dd}{\,\mathrm{d}}
\begin{document}
\section*{UnsolvedMath ID: 11000012}
\textbf{Source problem code:} catalogue entry
\subsection{Definitions and mathematical statement}
Let \(\Sigma_h\) be an oriented genus-\(h\) surface and \(\operatorname{Mod}_g\) the mapping class group of a genus-\(g\) surface. Given
\[
 \rho:\pi_1(\Sigma_h)\longrightarrow\operatorname{Mod}_g,
\]
Problem 2.6 asks for an algorithm or group-theoretic invariant deciding/detecting whether \(\rho\) has a holomorphic representative. Conventionally this means a map from a Riemann-surface structure on \(\Sigma_h\) to the genus-\(g\) moduli orbifold \(\mathcal M_g\), inducing \(\rho\) in the appropriate orbifold sense. The source extract does not state whether the base complex structure is fixed or variable, the genus range, or whether constant/isotrivial representatives are excluded.

The source says irreducibility is necessary but not sufficient, defining its condition literally as absence of an essential simple-closed-curve isotopy class \(\alpha\) fixed by every image element:
\[
 \nexists\alpha\quad
 \rho(\gamma)(\alpha)=\alpha\quad\text{for all }\gamma\in\pi_1(\Sigma_h).
\]
It expresses the belief that this condition is algorithmically decidable. This wording does not explicitly exclude an invariant multicurve whose components are permuted. Also, a constant holomorphic map has trivial, hence curve-fixing, monodromy, so the necessity statement needs omitted nonconstant or other source hypotheses. The reported necessity is not adopted unconditionally.

The surrounding program asks to construct and classify such representations up to conjugacy and compute invariants including
\[
 \nu(\rho)=\int_{\Sigma_h} f^*\omega_{\rm WP},
\]
where \(\omega_{\rm WP}\) is the Weil--Petersson \(2\)-form and \(f\) is any suitable map with \(f_*=\rho\). In a classifying-orbifold convention this can be expressed as \(\langle\rho^*[\omega_{\rm WP}],[\Sigma_h]\rangle\), provided the cohomology class and normalization are fixed. Using the coarse moduli space instead would not automatically identify its fundamental group with \(\operatorname{Mod}_g\); the map and well-definedness conventions require review.

The source links this computation to signatures of surface bundles over surfaces and the Gromov co-norm of \([\omega_{\rm WP}]\), without specifying that norm or a proposed bound. It regards unrestricted classification as unmanageable, noting surface-group surjections to free groups, and suggests first restricting to injective \(\rho\). It reports many irreducible injective examples via a Crisp--Wiest technique and Crisp--Farb [CF]. These are contextual reports, not an undecidability theorem or a completed classification.
\subsection{Short English statement}
Detect which surface-group representations into a mapping class group are realized holomorphically, then study their conjugacy classes and Weil--Petersson characteristic number. The source suggests starting with injective representations. Base-complex-structure, orbifold and nonconstant-map conventions are essential missing qualifications.
\subsection{Sources}
\begin{itemize}
\item Local UnsolvedMath registry, permanent ID \(11000012\), exact source code \texttt{catalogue entry}, Problem 2.6, ``Holomorphic representatives,'' PDF page 25. No individual proposer or proposed year is supplied.
\item Farb, edited volume \emph{Problems on Mapping Class Groups and Related Topics}: \url{https://www.math.uchicago.edu/~farb/papers/mcgbook.pdf}. The extract includes chapter/page furniture ``Some problems on mapping class groups and moduli space,'' page 19. The volume has 329 numbered theorem-style items with repeated local numbering distinguished by PDF page.
\item Supplied internal reference [CF] is described as J. Crisp with the chapter author; the research summary calls it Crisp--Farb and mentions a Crisp--Wiest technique. Full bibliography entries are absent. The summary dates the broader source discussion to Farb (2006).
\item Extraction pdf-font-structure, accessed 2026-07-29. Source evidence NEEDS\_REVIEW presents an open problem and requires a current primary-literature audit. Supplied OPEN-TRIAGE summary reports no located algorithm/invariant, unresolved classification and \(\nu\)-computations, and largely open co-norm/signature questions; no separate modern primary references are supplied.
\item Dataset attribution: Ulam AI / UnsolvedMath Contributors, CC BY 4.0. Source rights NEEDS\_REVIEW: publicly accessible author/publisher copy, redistribution terms requiring review.
\end{itemize}
\end{document}
